The project includes a large number of negative or mixed-signal probes.
These do not belong in the main narrative, but they are important for interpreting the paper's boundaries.

\paragraph{Value-head sweeps.}
We tested richer value heads including pairwise, interaction, heteroscedastic, low-rank, conditional-latent, and calibration-wrapper variants.
Taken together, these sweeps support a narrower conclusion than ``we only need a stronger head.''
Richer parameterizations often improve exact-state performance, but deployable predicted-state performance is much harder to improve robustly.
Many apparent gains are domain-specific, and many do not survive cross-domain comparison or repeatability checks.

\paragraph{Continuation outcome statistics.}
The richer structured-controller features eventually included quantities such as wrong-rate, token cost, and utility dispersion.
These were empirically useful, but they also moved the representation away from the original narrow success-distribution interpretation of \(S_t\).
This is why the main paper now uses the more conservative term \emph{continuation outcome statistics}.

\paragraph{Why these negative results matter.}
The point of this appendix is not to advertise more controller families.
It is to document that the final paper boundaries were chosen only after substantial unsuccessful attempts to eliminate the deployment gap by local architectural search alone.
Those failures strengthen the main interpretation of the paper:
the central obstacle is not merely missing head capacity, but reliable state identification and deployment under noise.

\paragraph{State-identification rescue probes.}
We also tested more targeted rescue probes aimed directly at the deployment gap rather than at generic head capacity.
These included high-determinacy-only training slices, pairwise/soft teacher targets, exact-state teacher distillation, and combinations of the best cleanup and teacher variants.
The resulting picture was again conditional rather than universal.
Some probes improved linear equations, others improved arithmetic, and some combinations reduced revision harm without improving regret.
But the gains did not stack into a stable rescue, and the best predicted-state variants still remained above the strongest learned one-dimensional baseline on both domains.
This is the key reason the main paper stops at a benchmark / mechanism / deployment-diagnosis claim.
